\section{Filters}
\subsection{Fundamentals}

\subsubsection{Basic Filters}
We give here the processing stage after the pre-amplifier:
\begin{figure}[H]
    \centering
    \includegraphics[width=0.5\linewidth]{Post_Pre_AMp_Processing.png}
    \caption{Processing stage after the Preamp.}
\end{figure}
Filters can either be passive or active:


\begin{tabular}{|p{8cm}|p{8cm}|}
\hline
\textbf{Passive} & \textbf{Active} \\
\hline
Only passive elements (R, L, C, transformers, gyrators, …) 
& Contain both passive and active components (transistors) \\
\hline
No need of power supply & Need of power supply \\
\hline
There is dissipation $P_{out} < P_{in}$ & $P_{out} \geq P_{in}$ \\
\hline
Very difficult to design, we cannot control the input/output impedances. 
Sources/load can modify the filter response. 
& We can decouple source/load impedance \\
\hline
\end{tabular}

\begin{figure}[H]
    \centering
    \includegraphics[width=0.9\linewidth]{Filters_TF.png}
    \caption{Transfer functions of basic filters.}
\end{figure}

The order of a filter is related to the slope of the cuts (e.g. double poles and zeros).
\begin{figure}[H]
    \centering
    \includegraphics[width=0.4\linewidth]{Filter_Order_Illustration.png}
    \caption{Relationship between slope's order and filter's order.}
\end{figure}

The ideal LPF has gain 1 before $f_c$, and 0 after (which means $-\infty \ \text{dB}$), while the phase shifts (which is not desirable in practice).

\begin{figure}[H]
    \centering
    \includegraphics[width=0.5\linewidth]{Ideal_LPF.png}
    \caption{Transfer function of the ideal LPF.}
\end{figure}

Real filters differ from the ideal one and we have many different types. We provide some examples below: 
\begin{figure}[H]
    \centering
    \includegraphics[width=0.6\linewidth]{Digital_LPF.png}
    \caption{Butterworth and Chebicev filters.}
\end{figure}

\textbf{The usual expression of transfer functions of filters is:}

\begin{itemize}
    \item Low-pass (LPF): 
    \[
    TF_{LP} = \frac{K}{D(s)} \quad \text{for the 1st order}, \quad 
    TF_{LP} = \frac{K}{D_2(s)} = \frac{K}{s^2 + b_1 s + b_0} \quad \text{for the 2nd order}
    \]
    \item High-pass (HPF):
    \[
    TF_{HP} = K \cdot \frac{N(s)}{D(s)}, \qquad 
    TF_{HP} = \frac{K s^2}{D_2(s)} \quad \text{for the 2nd order with two zeros at the origin}
    \]
    \item Band-pass (BPF):
    \[
    TF_{BP} = K \cdot \frac{N(s)}{D(s)}, \qquad 
    TF_{BP} = K \cdot \frac{s}{D_2(s)} \quad \text{for the 2nd order with one zero at the origin}
    \]
    \begin{itemize}
        \item For band-pass filters, we usually have the order of $N$ being lower than the order of $D$.
    \end{itemize}
\end{itemize}

The most generic 2\textsuperscript{nd} order transfer function is: 
\[
TF_2(s) = \frac{N_2(s)}{D_2(s)} = \frac{a_2 s^2 + a_1 s + a_0}{s^2 + b_1 s + b_0}
\]
which is referred to as \textbf{Biquadratic (Biquads) functions}.

The frequency response of a filter can be implemented in different circuits, talking about circuit topology: Sallen-Key biquad, Tow-Thomas biquad, Stable variable filter...  
They have different implementations with different sensitivity to parameters, they thus have different noise performance and complexity.

\clearpage 

\subsubsection{Sallen-Key Filter}
The LPF has a negative feedback for gain only while this filter also regards the positive feedback.  

In order to achieve stability, we attenuate the treble instead of boosting bass $G < 3$.  

We have:
\[
TF(s) = \frac{G}{s^2 \tau^2 + s(3-G)\tau + 1}, 
\qquad \tau = RC
\]

\begin{figure}[H]
    \centering
    % Left side: explanation + equations
    \begin{minipage}[c]{0.4 \linewidth}
        \centering
        \includegraphics[width=\linewidth]{Sallen_Key.png}
    \caption{Sallen-Key filter.}
    \end{minipage}\hfill
    \begin{minipage}[c]{0.45\linewidth}
        \centering
        \includegraphics[width=\linewidth]{Sallen_Key_dual.png}
    \caption{Sallen-Key filter dual representation.}
    \end{minipage}
\end{figure}

\subsubsection{Filter’s functionality}
All filters make variations on the signal’s spectrum, but for different reasons:
\begin{itemize}
    \item \textbf{Tone Control:} specialized in low/mid/treble boost/cut (see tube screamer).
    \item \textbf{Crossover:} frequency separation, each band is intended for different use cases.
    \item \textbf{Equalizer:} series of filters centred at different frequencies. There are graphic and parametric equalizers.
\end{itemize}

\begin{figure}[H]
    \centering
    \includegraphics[width=0.7\linewidth]{Filter_Functionality.png}
    \caption{Illustrtion of the frequency separation.}
    \label{fig:placeholder}
\end{figure}

\subsection*{Parametric Filters}
Let’s start from a LPF configuration and understand how to change the cutoff frequency without changing the in-band gain $G_{in,band}$:

\begin{figure}[H]
    \centering
    % Left side: explanation + equations
    \begin{minipage}[c]{0.6 \linewidth}
    We have $v^+ = v^- = 0$ from the ideal OPAMP equations, thus: 
    \begin{align*}
        & v^- = \frac{\frac{v_{in}}{R_1} + \frac{v_{out}}{R_2//C}}{\frac{1}{R_1} + \frac{1}{R_2//C}}=0
    \quad \Rightarrow \quad
    \frac{v_{in}}{R_1} + \frac{v_{out}}{R_2//C} = 0 \\
    \Rightarrow \quad & \frac{v_{out}}{v_{in}} = -\frac{\frac{R_2}{1+sCR_2}}{R_1}
    \quad \Rightarrow \quad
    \frac{v_{out}}{v_{in}} = -\frac{R_2}{R_1} \cdot \frac{1}{1+sC R_2} \\
    \Rightarrow \quad &
    \frac{v_{out}}{v_{in}} = -G_{in,band} \cdot \frac{1}{1+s\tau} , 
    \qquad G_{in,band} = \frac{R_2}{R_1} , \quad \tau = C R_2
    \end{align*}
    \end{minipage}\hfill
    \begin{minipage}[c]{0.35\linewidth}
        \centering
        \includegraphics[width=\linewidth]{Parametric_LPF.png}
    \caption{Parametric LPF scheme.}
    \end{minipage}
\end{figure}

\newpage
We can modify the cut-off frequency $ f_p = \tfrac{1}{2 \pi R_2 C}$ by changing:
\begin{itemize}
    \item $C$: very difficult (no potentiometer)
    \item $R_2$: easy, but $G_{in,band}$ will also change since $ GBWP = G_{in,band} \cdot f_p$ is constant
\end{itemize}

\begin{figure}[H]
    \centering
    \includegraphics[width=0.4\linewidth]{Parametric_LPF_TF.png}
    \caption{Transfer function of the parametric LPF.}
\end{figure}
In order not to change $G_{in,band}$, we may counterbalance ${R}_2$ increase with $R_1$ (dual potentiometer).  

\begin{figure}[H]
    \centering
    % Left side: explanation + equations
    \begin{minipage}[c]{0.6 \linewidth}
    We have $v^+ = v^- = 0$ from the ideal OPAMP equations, thus: 
    \begin{align*}
        & v^- = \frac{\tfrac{v_{in}}{P} + \tfrac{v_{out}}{P//C}}{\tfrac{1}{P} + \tfrac{1}{P//C}} = 0
\quad \Rightarrow \quad 
\tfrac{v_{in}}{P} + \tfrac{v_{out}}{P//C}=0 \\
\Rightarrow \quad & v_{out} = -\frac{v_{in}}{P}\cdot P//C
\quad \Rightarrow \quad \frac{v_{out}}{v_{in}} = -\frac{\tfrac{P/sC}{P+1/sC}}{P}=-\frac{\tfrac{P}{1+sCP}}{P} \\
\Rightarrow \quad & \frac{v_{out}}{v_{in}} =-\frac{P}{P} \frac{1}{1+sCP}
\quad \Rightarrow \quad \frac{v_{out}}{v_{in}}= -\frac{1}{1+sCP}
    \end{align*}
    \end{minipage}\hfill
    \begin{minipage}[c]{0.35\linewidth}
        \centering
        \includegraphics[width=\linewidth]{Parametric_LPF_dual.png}
    \caption{Parametric LPF with the dual potentiometer $P$.}
    \end{minipage}
\end{figure}
On the transfer function below, we do observe that $G_{in,band}$ is conserved even when modifying $P$.
\begin{figure}[H]
    \centering
    \includegraphics[width=0.4\linewidth]{Parametric_LPF_Dual_TF.png}
    \caption{Transfer function of the Parametric LPF with dual potentiometer $P$.}
\end{figure}

\begin{figure}[H]
    \centering
    % Left side: explanation + equations
    \begin{minipage}[c]{0.6 \linewidth}
    The parametric HPF has a similar design, we have here: 
\begin{align*}
& v^- = \frac{\tfrac{v_{in}}{P+C} + \tfrac{v_{out}}{P}}{\tfrac{1}{P+C} + \tfrac{1}{P}} = 0
\quad \Rightarrow \quad 
\tfrac{v_{in}}{P+C} + \tfrac{v_{out}}{P}=0 \\
\Rightarrow \quad & v_{out} = -\frac{v_{in}}{P+C}\cdot P
\quad \Rightarrow \quad 
\frac{v_{out}}{v_{in}} = -\frac{P}{P+1/sC} \\
\Rightarrow \quad & \frac{v_{out}}{v_{in}} =-\frac{scP}{1+sCP} 
\end{align*}
    \end{minipage}\hfill
    \begin{minipage}[c]{0.35\linewidth}
        \centering
        \includegraphics[width=\linewidth]{Parametric_HPF_Dual.png}
    \caption{Parametric HPF with the dual potentiometer $P$.}
    \end{minipage}
\end{figure}


Similarly to the parametric LPF, we see on the transfer function below that $G_{in,band}$ is conserved even when modifying $P$.
\begin{figure}[H]
    \centering
    \includegraphics[width=0.4\linewidth]{Parametric_HPF__TF.png}
    \caption{Transfer function of the parametric HPF with dual potentiometer $P$.}
\end{figure}

\subsection{Tone Control Filters}
Tone control filters are sometimes associated with equalization. More precisely, bass/treble control.  

\begin{figure}[H]
    \centering
    \includegraphics[width=0.7\linewidth]{Tone_Control_Filters_Gen.png}
    \caption{Audio equalization performed by tone control filters.}
\end{figure}

Let's start from a simple circuit as shown below: 
\begin{figure}[H]
    \centering
    \includegraphics[width=0.5\linewidth]{Simple circuit.png}
    \caption{Basic HPF scheme.}
\end{figure}

We want the transfer function to have more zeros in $f_z \neq 0$, and add a pole $f_p > f_z$ in order to obtain a low-pass shelf.  
We show below two different ways for achieving such a goal.

\begin{figure}[H]
    \centering
    % Left side: explanation + equations
    \begin{minipage}[c]{0.45 \linewidth}
    We have $v^+ = v^- = 0$ from the ideal OPAMP equations, thus: 
\begin{align*}
& v^- = \frac{\tfrac{v_{in}}{R//C} + \tfrac{v_{out}}{R_1}}{\tfrac{1}{R//C} + \tfrac{1}{R_1}} = 0
\quad \Rightarrow \quad 
\tfrac{v_{in}}{R//C} + \tfrac{v_{out}}{R_1}=0 \\
\Rightarrow \quad 
& v_{out} = -\frac{v_{in}}{R//C}\cdot R_1
\quad \Rightarrow \quad 
\frac{v_{out}}{v_{in}} = -\frac{R_1}{\tfrac{R/sC}{R+1/sC}} \\
\Rightarrow \quad & \frac{v_{out}}{v_{in}} = - \frac{R_1}{\tfrac{R}{1+sCR}} 
\quad \Rightarrow \quad \frac{v_{out}}{v_{in}}=-\frac{R_1}{R} \cdot \left( 1+sCR\right)
\end{align*}
    \end{minipage}\hfill
    \begin{minipage}[c]{0.45\linewidth}
        \centering
        \includegraphics[width=\linewidth]{LPF_Shelf_1.png}\caption{Introducing a high pass shelf with a parallel resistor $R$ on the input.}
    \end{minipage}
\end{figure}


Expanding the circuit's structure, we can introduce a highpass shelf by adding a series resistor on the input: 
\begin{figure}[H]
    \centering
    \includegraphics[width=0.6\linewidth]{Tone_Control_2.png}
    \caption{Corrected configuration of the bass cut control filter.}
\end{figure}

Similarly to what we did before, we observe : 
\begin{align*}
& \frac{v_{in}}{R//C+R_1} + \frac{v_{out}}{R_1}=0 
\quad \Rightarrow \quad
v_{out} = -\frac{v_{in}}{R//C+R_1}\cdot R_1
\quad \Rightarrow \quad 
\frac{v_{out}}{v_{in}} = -\frac{R_1}{R_1 +  \tfrac{R/sC}{R+1/sC}} \\
\Rightarrow \quad
& 
\frac{v_{out}}{v_{in}} = - \frac{R_1}{R_1 + \tfrac{R}{1+sCR}} 
\quad \Rightarrow \quad 
\frac{v_{out}}{v_{in}}=-\frac{R_1 \left(1+sCR \right)}{R + R_1 + sCRR_1} 
\quad \Rightarrow \quad 
\frac{v_{out}}{v_{in}} = - \frac{R_1}{R+R_1}\frac{ 1+sCR }{1 + sC\left(R // R_1 \right)} 
\end{align*}

We want to control the bass-cut amount only by changing $R$.  
The bass-cut gain is
\[
G_{BC} = \frac{v_{out}}{v_{in}} 
= -\frac{Z_F}{Z_{IN}} 
= - \frac{R_1}{R+R_1} \cdot \frac{1+s\tau_z}{1+s\tau_p}
\]
Where:
\begin{itemize}
    \item $\dfrac{R_1}{R+R_1} = G_{BC}(0)$ the static gain
    \item $\tau_z$: the zero, related to $R \cdot C$
    \item $\tau_p$: the pole, related to $ C \left( R//R_1 \right)$
\end{itemize}


Intuitively, we can switch $Z_F \leftrightarrow Z_{in}$ to implement a bass-boost:
\[
G_{BB} = \frac{1}{G_C} = \frac{R+R_1}{R_1} \cdot \frac{1+s\tau_p}{1+s\tau_z},
\qquad
\tau_p = C \left( R//R_1 \right), \quad \tau_z=RC
\]
\begin{figure}[H]
    \centering
    \includegraphics[width=0.7\linewidth]{Tone_Control_Bass_Boost.png}
    \caption{Scheme of the bass boost.}
\end{figure}


The overall bass-control configuration is:
\[
\text{Bass Control: } 
\quad \max \, \text{boost} = +K\,\text{dB}, 
\quad \max \, \text{cut} = -K\,\text{dB},
\quad \text{mid point}= 0 \ \text{dB}
\quad f_z = \frac{1}{2 \pi R \cdot C}, 
\]
\begin{figure}[H]
    \centering
    \includegraphics[width=0.5\linewidth]{Overall_Bass_Control.png}
    \caption{Complete scheme of the bass-control configuration.}
\end{figure}

Similarly, the overall treble-control configuration is:
\[
\text{Treble Control: } 
\quad \max \, \text{boost} = +K\,\text{dB}, 
\quad \max \, \text{cut} = -K\ ,\text{dB}, 
\quad \text{mid point} = 0 \ \text{dB}
\quad f_p = \frac{1}{2 \pi R \cdot C}, 
\]
\begin{figure}[H]
    \centering
    \includegraphics[width=0.45\linewidth]{Overall_Treble_Control.png}
    \caption{Complete scheme of the treble-control configuration.}
\end{figure}

Finally, the Baxandall tone control merges the previous bass/treble controls into a single stage:
\begin{figure}[H]
    \centering
    \includegraphics[width=0.8\linewidth]{Bakandall_Tone_Control.png}
    \caption{Scheme of the Baxandall tone control.}
\end{figure}



\subsection{Notes from the slides}

\paragraph{Active Filters and Tone Control Circuits: }

Filters are essential in both analog and digital audio systems.  
They shape the frequency content of signals, attenuating or amplifying specific bands to achieve tonal balance, noise reduction, or crossover between drivers.

In this part we discuss:
\begin{itemize}
  \item Tone control circuits (Bass and Treble),
  \item Active filter topologies such as the Sallen–Key configuration,
  \item Crossover and equalization networks in professional audio.
\end{itemize}

\paragraph{Tone Control Filters: }
Tone controls allow independent adjustment of low and high frequencies.  
A typical \textbf{bass} control boosts or cuts frequencies below about $100\,\text{Hz}$–$300\,\text{Hz}$,  
while a \textbf{treble} control acts above $5\,\text{kHz}$–$10\,\text{kHz}$.

\paragraph{The Baxandall Circuit: }
One of the most widely used tone control circuits is the \textbf{Baxandall network}, introduced by Peter Baxandall in 1952.  
It employs active feedback to achieve smooth, symmetrical boost and cut characteristics with minimal distortion.


\textbf{Example: Luxman L-215 Amplifier: }
In high-fidelity amplifiers such as the \emph{Luxman L-215}, the tone control stage based on the Baxandall configuration is inserted between the preamplifier and the main amplifier stages.  
This placement ensures that tonal corrections are applied before the power amplification, avoiding overload or distortion.

\paragraph{Filter Order and Attenuation: }
Filters are classified by their order, corresponding to the highest power of \(s\) in the denominator of their transfer function.  
Each additional order increases the roll-off steepness:

\begin{center}
\begin{tabular}{|c|c|c|}
\hline
\textbf{Order} & \textbf{Attenuation per Octave} & \textbf{Attenuation per Decade} \\
\hline
1 & $-6$ dB/oct & $-20$ dB/dec \\
2 & $-12$ dB/oct & $-40$ dB/dec \\
3 & $-18$ dB/oct & $-60$ dB/dec \\
\hline
\end{tabular}
\end{center}

The most common responses are:
\begin{itemize}
  \item \textbf{Butterworth} – maximally flat amplitude response.
  \item \textbf{Chebyshev} – equiripple in the passband, steeper roll-off.
\end{itemize}

\begin{figure}[H]
    \centering
    \includegraphics[width=0.5\linewidth]{Order_Filter.png}
    \caption{Illustration of the different attenuations per octave for different orders of filters.}
\end{figure}

\paragraph{Active Filters: the Sallen–Key Topology}
Active filters use operational amplifiers to implement low-, high-, or band-pass responses without inductors.  
The \textbf{Sallen–Key topology} is among the simplest and most robust second-order filter designs.

\textbf{Low-Pass Filter:}
\[
TF_{LP}(s) = \frac{G}{s^2 R_1 R_2 C_1 C_2 + s \left[(3 - G) R_2 C_2 + R_1 C_1 \right] + 1}.
\]

\textbf{High-Pass Filter:}
\[
TF_{HP}(s) = \frac{G s^2 R_1 R_2 C_1 C_2}{s^2 R_1 R_2 C_1 C_2 + s \left[(3 - G) R_1 C_1 + R_2 C_2 \right] + 1}.
\]
These filters provide a slope of \(-12\,\text{dB/oct}\) (\(-40\,\text{dB/dec}\)) and are widely used in professional mixing consoles.
\begin{figure}[H]
    \centering
    \includegraphics[width=0.5\linewidth]{Sallen_Key_Slides.png}
    \caption{Sallen–Key low-pass and high-pass configurations used for tone control and crossover applications.}
\end{figure}

\paragraph{Crossover Filters: }
\textbf{Crossover filters} divide the audio spectrum into separate bands—low, mid, and high—each feeding a dedicated amplifier and loudspeaker driver:
\[
\begin{aligned}
f_\text{L--M} &\approx 200\,\text{Hz}, \qquad
f_\text{M--H} &\approx 2\,\text{kHz}.
\end{aligned}
\]
\begin{figure}[H]
    \centering
    \includegraphics[width=0.5\linewidth]{Crossover_Filters.png}
    \caption{Illustration of the per-band separation of the spectrum.}
\end{figure}

Using active components ensures precise frequency division and better phase alignment compared to passive crossovers.

\paragraph{Graphic Equalizers: }
A \textbf{graphic equalizer} provides independent boost or cut on multiple frequency bands (typically 10, 15, or 31 bands).  
Each band corresponds to a band-pass filter centered at \(f_i\), controlled by a slider (fader).

When the slider is centered, the response is flat.  
Moving it upward \emph{boosts} the band around \(f_i\); downward movement \emph{cuts} it:
\[
v_{out} = v_{in}\,(1 \pm \beta G_{f_i}),
\quad \text{where \(G_{f_i}\) is the gain of the band-pass cell.}
\]
\begin{figure}[H]
    \centering
    \includegraphics[width=0.3\linewidth]{Graphic_EQ.png}
    \caption{Simple representation  of the graphic EQ control over frequencies $f_i$.}
\end{figure}

\paragraph{Effect of the Quality Factor \(Q\).}
The selectivity of each band depends on the quality factor:
\[
Q = \frac{f_0}{\Delta f} = \frac{f_0}{BW}.
\]
High-\(Q\) filters affect only narrow frequency ranges, while low-\(Q\) filters create smoother, overlapping corrections.
\begin{figure}[H]
    \centering
    \includegraphics[width=0.8\linewidth]{Q_Factor.png}
    \caption{Influence of the quality factor \(Q\) on the equalizer band response. Higher \(Q\) produces a narrower, more selective peak or dip. Linear scale on the left, log scale on the right.}
\end{figure}

\paragraph{Equalizer circuit: } We show below the computation of the output of the circuit for different configurations. \textbf{CHECK SLIDE L10 - Graphic equalizer details.pdf}

\begin{figure}[H]
    \centering
    % Left side: explanation + equations
    \begin{minipage}[c]{0.45 \linewidth}
        \centering
        \includegraphics[width=\linewidth]{EQ_circuit_1.png}
        \caption{Starting Configuration.}
    \end{minipage}\hfill
    \begin{minipage}[c]{0.45\linewidth}
        \centering
        \includegraphics[width=\linewidth]{EQ_Circuit_2.png}
        \caption{Flat response, cursor at the center.}
    \end{minipage}
    \vspace{0.5cm}
   \begin{minipage}[c] {0.45 \linewidth}
        \centering
        \includegraphics[width=\linewidth]{EQ_circuit_3.png}
        \caption{Boost around $f_1$, cursor at the top.}
    \end{minipage}\hfill
    \begin{minipage}[c] {0.45 \linewidth}
        \centering
        \includegraphics[width=\linewidth]{EQ_circuit_4.png}
    \caption{Cut around $f_1$, cursor at the bottom}
    \end{minipage}
\end{figure}


\paragraph{Parametric Equalizers}
Unlike graphic equalizers, \textbf{parametric equalizers} allow full control over:
\begin{itemize}
  \item \textbf{Center frequency} \(f_0\),
  \item \textbf{Gain} (boost or cut),
  \item \textbf{Bandwidth or quality factor \(Q\)}.
\end{itemize}

The core element is a tunable band-pass filter, typically implemented with operational amplifiers.

\begin{figure} [H]
    \centering
    \includegraphics[width=0.5\linewidth]{Param_EQ.png}
    \caption{Response of a parametric equalizer showing adjustable \(f_0\), gain, and \(Q\).  
This configuration enables precise shaping of the signal spectrum.}
\end{figure}

Digital implementations of parametric and graphic equalizers are now standard in DAWs and mixing consoles, using IIR or FIR filters for accurate and recallable tone shaping.

\textbf{Summary}
\begin{itemize}
  \item Tone controls (Bass/Treble) use active feedback networks such as Baxandall filters.
  \item Active filters (e.g., Sallen–Key) provide stable, low-distortion second-order responses for tone control and crossovers.
  \item Equalizers—graphic and parametric—extend filtering concepts to multiple or variable frequency bands, giving engineers fine spectral control.
  \item Digital realizations replicate these analog behaviors with high precision and flexibility.
\end{itemize}

